\section{Defensa ante contenido desconocido: Predictive Detection, calibración y revisión humana}

Tanto las herramientas generativas como el contenido recién subido producen objetos que el reference set no ha visto; por eso, la siguiente capa es la predictive detection. Lo esencial aquí no es «si el modelo puede asignar una puntuación», sino cómo convertir el score en una cola de revisión operable y sostenible: el umbral debe ajustarse según el riesgo y la capacidad, la salida debe pasar por calibration y la decisión final debe conservar la human review y los límites de apelación.

\subsection{Del classifier score al verified outcome}

Un \term{predictive classifier} (clasificador predictivo) aprende, a partir de los datos de entrenamiento, la relación estadística entre la entrada y las etiquetas de riesgo, y produce una clase o un score. La \term{human review} (revisión humana) consiste en que personal capacitado, sujeto a restricciones de permisos y a medidas de protección de su salud, tome decisiones con base en el contexto y la policy; no se trata de simplemente «echar un vistazo al resultado del modelo». La salida del modelo es una triage signal, no una conclusión jurídica.

\lecturefigure{04-predictive-classifier.png}{La Predictive AI amplía la cobertura más allá de los hash conocidos, pero debe pasar por la cola y la revisión.}{Subtítulos locales 10:10--13:40; mecanismo oficial de Thorn Safer Predict; redibujo conceptual.}

\begin{knowledgebox}{Lectura de la figura: el significado del score depende del entrenamiento y la calibración}
El modelo produce primero un raw score; después, la plataforma asigna cada objeto a un nivel de prioridad según el threshold; el reviewer ve el case context filtrado por permisos y escribe el verified outcome de vuelta en el sistema controlado. La magnitud del raw score no puede compararse directamente entre model versions, tipos de medio o entornos de cliente. Antes del despliegue debe precisarse si el score es un ranking, una estimated probability o un uncalibrated margin, y deben registrarse model/version y threshold/version.
\end{knowledgebox}

\teachervoice{Cordua usa la evolución de Safer Match a Safer Predict para explicar que, cuando los atacantes o los modelos generativos crean objetos nuevos sin cesar, ya no basta con buscar solo hash conocidos; pero el nombre «Predict» tampoco significa que el sistema pueda prescindir de las personas. El punto central de la clase es ampliar la cobertura de candidatos y adelantar los elementos que más vale la pena revisar.}

\subsection{Precision, recall, calibration y queue capacity}

Esta sección reúne en una misma figura las métricas del modelo y las restricciones operativas. La \term{precision} (precisión) indica qué proporción de los objetos que el sistema marca como positivos se verifica como positiva; el \term{recall} (exhaustividad) indica qué proporción de todos los positivos reales encuentra el sistema; la \term{calibration} (calibración) indica si la probabilidad predicha coincide con la frecuencia real observada a largo plazo. Ninguna de las tres puede optimizarse por separado sin considerar la base rate, el label process y la reviewer capacity.

\lecturefigure{10-evaluation-queue.png}{La calidad del modelo y la capacidad de revisión deben diseñarse juntas; el umbral es una decisión operativa, no una constante pura del modelo.}{Subtítulos locales 10:10--16:20; enfoque de medición de riesgos del NIST AI RMF; redibujo conceptual.}

\begin{importantbox}{Definición de las métricas y explicación de los símbolos}
Si se denotan los verdaderos positivos, falsos positivos y falsos negativos como $TP,FP,FN$, respectivamente, entonces
\[
\mathrm{precision}=\frac{TP}{TP+FP},\qquad
\mathrm{recall}=\frac{TP}{TP+FN}.
\]
La precision determina cuántos elementos inválidos encuentra el revisor entre las alertas, y el recall determina cuántos riesgos reales omite el sistema. Si el número diario de candidatos es $N$, la tasa de alertas con el umbral dado es $q(\tau)$, la capacidad diaria de cada revisor es $c$ y el número de revisores es $R$, una operación estable requiere al menos
\[
Nq(\tau) \lesssim Rc,
\]
donde $\tau$ es el umbral. Esta desigualdad es una restricción de capacidad; no significa que el objetivo de seguridad se reduzca a «evitar que la cola se acumule».
\end{importantbox}

\begin{warningbox}{Una accuracy alta puede ocultar una base rate extremadamente baja}
En un entorno donde los positivos reales son escasísimos, un modelo que casi siempre predice negativo puede tener una accuracy muy alta. La plataforma debe reportar precision/recall estratificados por tipo de medio, región, idioma, escenarios relacionados con la edad y punto de entrada del producto, y monitorear la distribution shift. Un cambio de umbral afecta los falsos negativos, pero también altera la exposición de los revisores y la carga de los reportes obligatorios por ley.
\end{warningbox}

\subsection{La moderator triage es parte del sistema de seguridad}

La revisión humana no es un último paso ilimitado ni gratuito. El \term{triage} (triaje/priorización) ordena los candidatos según gravedad, confianza, urgencia y posibilidad de actuar; las buenas herramientas, además, ocultan la información visual no esencial, agrupan los objetos duplicados, ofrecen el contexto mínimo necesario y dan soporte a la rotación de turnos y a mecanismos de salud mental. Al leer la figura de esta sección hay que observar a la vez la «calidad de la decisión» y la «exposición humana»; no basta con perseguir el rendimiento.

\lecturefigure{05-moderator-triage.png}{El triage debe maximizar el valor de protección y, al mismo tiempo, minimizar la exposición sensible innecesaria de los revisores.}{Subtítulos locales 13:40--16:20; redibujo conceptual basado en la discusión de Safer Review.}

\begin{knowledgebox}{Lectura de la figura: un case package es mejor que una cola de archivos sueltos}
La capa de entrada aplica primero hash, ordenamiento por modelo y agrupación de duplicados; el workbench muestra solo la evidencia mínima necesaria para aplicar la policy; la reviewer decision y su rationale se registran de forma estructurada; después, el case package incorpora la account, el tiempo, los eventos relacionados y el action history. Así se reducen las visualizaciones repetidas y, además, la segunda revisión, el muestreo de calidad, los reportes y la incident review pueden realizarse sobre un mismo objeto auditable.
\end{knowledgebox}

\begin{warningbox}{El reviewer wellbeing no puede depender de la «fortaleza psicológica»}
La organización necesita exposure minimization, capacitación, rotación, descansos, apoyo profesional, monitoreo de accesos y clear escalation. Los falsos positivos del modelo tienen un costo humano real; por eso, un false positive no es solo una cifra en el dashboard. A la inversa, tampoco se puede ajustar el threshold hasta provocar omisiones sistemáticas con tal de reducir la exposición. Las políticas de producto, de modelo y de personal deben asumir el riesgo en conjunto.
\end{warningbox}

\subsection{Resumen de la sección}

El criterio de éxito de la predictive detection no es la accuracy en un leaderboard, sino que una calibrated signal ayude a los revisores a atender antes los objetos de alto riesgo dentro de una cola sostenible. Las métricas del modelo, el umbral, el case packaging, la salud del personal y los mecanismos de apelación deben validarse en conjunto.
